\newcounter{english}
\documentclass{article}

% packages
\usepackage{amsmath, amsthm, thmtools, amsfonts, amssymb, luacode, catchfile, tikzducks, hyperref, ifthen}
\ifcsname c@kobocompile\endcsname
	\usepackage[a5paper, total={1072pt, 1448pt}, margin=10pt, includeheadfoot]{geometry} % set page margins
\else
	\usepackage[a4paper, margin=50pt, includeheadfoot]{geometry}
\fi
\usepackage[shortlabels]{enumitem}
\usepackage[skip=3pt, indent=0pt]{parskip}

% language
\usepackage[bidi=basic, layout=tabular, provide=*]{babel}
\ifcsname c@english\endcsname
	\babelprovide[main, import]{english}
\else
	\babelprovide[main, import]{hebrew}
	\babelprovide{rl}
\fi
%\babelfont{rm}{Libertinus Serif}
\babelfont{rm}[Renderer=Harfbuzz]{Libertinus Serif}
\babelfont{sf}{Libertinus Sans}
\babelfont{tt}{Libertinus Mono}

% style
\AddToHook{cmd/section/before}{\clearpage}	% Add line break before section
\linespread{1.3}
\setcounter{secnumdepth}{0}		% Remove default number tags from sections, this won't do well with theorems
\AtBeginDocument{\setlength{\belowdisplayskip}{3pt}}
\AtBeginDocument{\setlength{\abovedisplayskip}{3pt}}
\graphicspath{ {../images/} }

% operators
\DeclareMathOperator\cis{cis}
\DeclareMathOperator\Sp{Sp}
\DeclareMathOperator\tr{tr}
\DeclareMathOperator\im{Im}
\DeclareMathOperator\re{Re}
\DeclareMathOperator\diag{diag}
\DeclareMathOperator*\lowlim{\underline{lim}}
\DeclareMathOperator*\uplim{\overline{lim}}
\DeclareMathOperator\rng{rng}
\DeclareMathOperator\Sym{Sym}
\DeclareMathOperator\Arg{Arg}
\DeclareMathOperator\Log{Log}
\DeclareMathOperator\dom{dom}
\DeclareMathOperator\supp{Supp}
\DeclareMathOperator\var{Var}
\DeclareMathOperator\cov{Cov}

% commands
%\renewcommand\qedsymbol{\textbf{מש''ל}}
%\renewcommand\qedsymbol{\fbox{\emoji{lizard}}}
\newcommand{\Aa}[0]{\mathcal{A}}
\newcommand{\Bb}[0]{\mathcal{B}}
\newcommand{\CC}[0]{\mathbb{C}}
\newcommand{\Cc}[0]{\mathcal{C}}
\newcommand{\EE}[0]{\mathbb{E}}
\newcommand{\FF}[0]{\mathbb{F}}
\newcommand{\Ff}[0]{\mathcal{F}}
\newcommand{\Ii}[0]{\mathcal{I}}
\newcommand{\Gg}[0]{\mathcal{G}}
\newcommand{\Ll}[0]{\mathcal{L}}
\newcommand{\Mm}[0]{\mathcal{M}}
\newcommand{\NN}[0]{\mathbb{N}}
\newcommand{\Nn}[0]{\mathcal{N}}
\newcommand{\PP}[0]{\mathbb{P}}
\newcommand{\Pp}[0]{\mathcal{P}}
\newcommand{\QQ}[0]{\mathbb{Q}}
\newcommand{\RR}[0]{\mathbb{R}}
\newcommand{\Rr}[0]{\mathcal{R}}
\newcommand{\Ss}[0]{\mathcal{S}}
\newcommand{\TT}[0]{\mathbb{T}}
\newcommand{\Uu}[0]{\mathcal{U}}
\newcommand{\Vv}[0]{\mathcal{V}}
\newcommand{\Ww}[0]{\mathcal{W}}
\newcommand{\ZZ}[0]{\mathbb{Z}}
\newcommand{\acts}[0]{\circlearrowright}
\newcommand{\explain}[2] {
	\begin{flalign*}
		 && \text{#2} && \text{#1}
	\end{flalign*}
}
\newcommand{\maketitleprint}[0]{ \begin{center}
	%\begin{tikzpicture}[scale=3]
	%	\duck[graduate=gray!20!black, tassel=red!70!black]
	%\end{tikzpicture}	
	\includegraphics[width=6cm]{cover}
\end{center}
}

% theorem commands
\newtheoremstyle{c_remark}
	{}	% Space above
	{}	% Space below
	{}% Body font
	{}	% Indent amount
	{\bfseries}	% Theorem head font
	{}	% Punctuation after theorem head
	{.5em}	% Space after theorem head
	{\thmname{#1}\thmnumber{ #2}\thmnote{ \normalfont{\text{(#3)}}}}	% head content
\newtheoremstyle{c_definition}
	{3pt}	% Space above
	{3pt}	% Space below
	{}% Body font
	{}	% Indent amount
	{\bfseries}	% Theorem head font
	{}	% Punctuation after theorem head
	{.5em}	% Space after theorem head
	{\thmname{#1}\thmnumber{ #2}\thmnote{ \normalfont{\text{(#3)}}}}	% head content
\newtheoremstyle{c_plain}
	{3pt}	% Space above
	{3pt}	% Space below
	{\itshape}% Body font
	{}	% Indent amount
	{\bfseries}	% Theorem head font
	{}	% Punctuation after theorem head
	{.5em}	% Space after theorem head
	{\thmname{#1}\thmnumber{ #2}\thmnote{ \text{(#3)}}}	% head content

\ifcsname c@english\endcsname
	\theoremstyle{plain}
	\newtheorem{theorem}{Theorem}[section]
	\newtheorem{lemma}[theorem]{Lemma}
	\newtheorem{proposition}[theorem]{Proposition}
	\newtheorem*{proposition*}{Proposition}
	%\newtheorem{corollary}[theorem]{אין חלופה עברית}

	\theoremstyle{definition}
	\newtheorem{definition}[theorem]{Definition}
	\newtheorem*{definition*}{Definition}
	\newtheorem{example}{Example}[section]
	\newtheorem{exercise}{Exercise}[section]

	\theoremstyle{remark}
	\newtheorem*{remark}{Remark}
	\newtheorem*{solution}{Solution}
	\newtheorem{conclusion}[theorem]{Conclusion}
	\newtheorem{notation}[theorem]{Notation}
\else
	\theoremstyle{c_plain}
	\newtheorem{theorem}{משפט}[section]
	\newtheorem{lemma}[theorem]{למה}
	\newtheorem{proposition}[theorem]{טענה}
	\newtheorem*{proposition*}{טענה}
	%\newtheorem{corollary}[theorem]{אין חלופה עברית}

	\theoremstyle{c_definition}
	\newtheorem{definition}[theorem]{הגדרה}
	\newtheorem*{definition*}{הגדרה}
	\newtheorem{example}{דוגמה}[section]
	\newtheorem{exercise}{תרגיל}[section]

	\theoremstyle{c_remark}
	\newtheorem*{remark}{הערה}
	\newtheorem*{solution}{פתרון}
	\newtheorem{conclusion}[theorem]{מסקנה}
	\newtheorem{notation}[theorem]{סימון}
\fi

% Questions related commands
\newcounter{question}
\setcounter{question}{1}
\newcounter{sub_question}
\setcounter{sub_question}{1}

\ifcsname c@english\endcsname
	\newcommand{\question}[1][0]{
		\ifthenelse{#1 = 0}{}{\setcounter{question}{#1}}
		\section{Question \arabic{question}}
		\addtocounter{question}{1}
		\setcounter{sub_question}{1}
	}

	\newcommand{\subquestion}[1][0]{
		\ifthenelse{#1 = 0}{}{\setcounter{sub_question}{#1}}
		\subsection{Part \alph{sub_question}}
		\addtocounter{sub_question}{1}
	}
\else
	\newcommand{\question}[1][0]{
		\ifthenelse{#1 = 0}{}{\setcounter{question}{#1}}
		\section{שאלה \arabic{question}}
		\addtocounter{question}{1}
		\setcounter{sub_question}{1}
	}

	\newcommand{\subquestion}[1][0]{
		\ifthenelse{#1 = 0}{}{\setcounter{sub_question}{#1}}
		\subsection{סעיף \localecounter{letters.gershayim}{sub_question}}
		\addtocounter{sub_question}{1}
	}
\fi

% import lua and start of document
\directlua{common = require ('../common')}

\GetEnv{AUTHOR}

% headers
\author{\AUTHOR}
\date\today

\title{Final Exercise Answer Sheet --- Logic Theory (2), 80424}

\DeclareMathOperator{\PA}{PA}
\DeclareMathOperator{\Coll}{Coll}
\DeclareMathOperator{\Ind}{Ind}
\DeclareMathOperator{\Sat}{Sat}

\begin{document}
\maketitle
\maketitleprint[yellow]

\question{}
\subquestion{}
Let $U \subseteq \Pp(\NN)$ be a non-principal ultrafilter, let $\langle \Mm_n \mid n < \omega \rangle$ be a sequence of $L$-structures, and $\Mm = \prod_{n < \omega} \Mm_n / U$. \\
We will show that for every countable consistent set of formulas $\Gamma(x)$ with parameters from $M$ is realized in $\Mm$, namely that $\Mm$ is countably saturated.
\begin{proof}
	Take a coverage $\langle \Sigma_n(x) \mid n < \omega \rangle \subseteq \Gamma(x)$ such that $|\Sigma_n| < \omega$ for all $n$.
	Then $\Sigma_n$ is realized, and let $[f_n] \in M$ be such that $\Mm \models \Sigma_n([f_n])$.
	Then $a_n = \{ j < \omega \mid \Mm_n \models \Sigma_n(f_n(j)) \} \in U$.
	Filters are closed to intersection, then let us assume that $a_{n + 1} \subseteq a_n$, otherwise we could define,
	\[
		g_{n + 1}(i) = \begin{cases}
			f_{n + 1}(i) & i \in a_n \\
			c_n & \text{otherwise}
		\end{cases}
	\]
	where $c_n \in M_n$ is some arbitrary value.

	We now take $a = \bigcap_{n < \omega} a_n$ and $[f] \in M$ such that for $n \in a$, $f(n) \in \{ f_i(n) \mid i < \omega \}$.
	If $a \in U$ then $\Mm \models \Gamma([f])$, then let us assume $a \notin U$, conversely $a^C = \NN \setminus a \in U$.
	$a^C \cap a_n \in U$ for all $n$ and therefore $a^C \cap a_n \ne \emptyset$.
	It immediately follows that $\emptyset \in U$, a contradiction.
\end{proof}

\subquestion{}
We define $\sigma$-complete ultrafilter $U$ as an ultrafilter such that it is closed to countable intersections. \\
Let $U$ be some $\sigma$-complete ultrafilter, $L = \{ = \}$, $\Mm = {(\NN, =)}^I / U$ for some index set $I$. \\
We will show that $|M| = \omega$ and deduce that $\Mm$ is not countably saturated.
\begin{proof}
	Directly by Łoś theorem and sentence of the form $\varphi_N = \bigwedge_{n < N} \exists x (x \ne c_n)$ we deduce that $|M| \ge \omega$.
	Define $C_x = {\{ x \}}^I$ the constant function, we will show that for every $[f] \in M$ there is $n < \omega$ such that $[f] = [C_n]$.
	Note that this is equivalent to the claim that $\{ j < \omega \mid f(j) = n \} \in U$.
	We will assume otherwise in contradiction, then $a_n = \{ j < \omega \mid f(j) \ne n \}$ is in $U$, and $a_n \cap a_m$ is non-empty for all $n \ne m$.
	We take $a = \bigcap_{n < \omega} a_n$, $U$ is $\sigma$-complete therefore $a \in U$.
	It follows that $f(j) \ne n$ for all $j \in I, n < \omega$, a contradiction to $f$ being $I \to \NN$ function.
\end{proof}

\subquestion{}
We will show that if $U$ is an ultrafilter on some indices set $I$ such that $U$ is not $\sigma$-complete and $\langle \Mm_i \mid i \in I \rangle$, \\
then $\Mm = \prod_{i \in I} \Mm_i / U$ is countably saturated.
\begin{proof}
	By $\sigma$-incompleteness we can assume that there is decreasing chain $\langle u_n \mid j < \omega \rangle \subseteq U$ such that $\bigcap_{n < \omega} u_n \notin U$.
	We can assume without loss of generality that $\bigcap u_n = \emptyset$ as $I \setminus \bigcap u_n \in U$ and therefore we can take $u_n \setminus \bigcap u_m$.

	Let $\Gamma(x)$ be countably realized set of formulas.
	By countability let us denote $\Gamma(x) = \{ \gamma_n \mid n < \omega \}$.
	For every $N < \omega$ we also define $\Gamma_N = \{ \gamma_n \mid n < N \}$.
	Then $\Gamma_N$ is finite set of formulas, and realized by $[f_N] \in M$.
	We can take $f_N$ such that,
	\[
		a_N = \{ j \in I \mid \Mm_j \models \Gamma_N(f_N(j)) \}
		\subseteq \{ j \in I \mid \Mm_j \models \Gamma_N(f_M(j)) \}
	\]
	for $M < N$ by finite intersections.
	The sequence $\langle a_n \mid n < \omega \rangle \subseteq U$ is decreasing.
	Let us take $U_n = u_n \cap a_n$ for every $n$, clearly $U_n \in U$.
	Note that for every $i \in I$ there is maximal $n < \omega$ such that $i \in U_n$, and let us define $f(i) = f_n(i)$.
	Indeed $f \in M$ and,
	\[
		\Mm \models \gamma_n([f])
		\iff \{ j \in I \mid \Mm_j \models \gamma_n(f(j)) \} \in U
		\impliedby \{ j \in I \mid \Mm_j \models \gamma_n(f_n(j)) \} \in U
	\]
	but the latter holds directly by the definition of $f$.
\end{proof}


\end{document}
